\documentclass{article}
\usepackage{graphicx}
\usepackage{amsmath,amssymb,amsthm}

\title{Quadratic convergence in gradient descent via a golden stepsize schedule}
\author{Boris Andrews}
\date{\today}

\newtheorem{proposition}{Proposition}
\newtheorem{conjecture}{Conjecture}

\begin{document}

\maketitle

\section{Introduction}

Gradient descent with a constant stepsize is known to be suboptimal on ill-conditioned
quadratic objectives. For
\[
f(x)=\frac{1}{2}x^\top A x-b^\top x,\qquad A=A^\top \succ 0,
\]
the error $e_k=x_k-x_\star$ evolves as
\[
e_{k+1}=(I-h_kA)e_k,\qquad
e_n=p_n(A)e_0,\qquad
p_n(t)=\prod_{k=1}^n(1-h_kt),\quad p_n(0)=1.
\]
Thus the convergence problem becomes polynomial design: choose $\{h_k\}$ so that
$\max_{\lambda\in[\lambda_{\min},\lambda_{\max}]}|p_n(\lambda)|$ is small.

Classical Chebyshev semi-iterative methods solve this minimax problem (up to scaling)
when spectral bounds are known. The roots of the corresponding polynomials follow an
arcsine-type distribution in the large-$n$ limit. This connects naturally to recent
interest in non-monotone and cyclic stepsize schedules, including schedules that
intentionally use occasional long steps.

This manuscript explores a specific deterministic schedule based on the golden ratio
\[
\phi=\frac{1+\sqrt5}{2},
\]
motivated by the low-discrepancy sequence generated by irrational rotation. The central
empirical observation is that placing polynomial roots at $\sin^2(\pi\phi k)$ yields
remarkably bounded oscillatory behavior and, in experiments on quadratics, an
accelerated decay consistent with $O(1/n^2)$ in function error.

\section{Golden stepsize schedule}

Let $0<a<b\le 1$. Define
\[
r_k=a+(b-a)\sin^2(\pi\phi k),\qquad k\ge 1,
\]
and use inverse stepsizes $r_k$ with scaling parameter $\eta>0$:
\[
h_k=\frac{\eta}{r_k}.
\]
Equivalently,
\[
h_k=\frac{\eta}{a+(b-a)\sin^2(\pi\phi k)}
=\frac{\eta}{\,b-(b-a)\cos^2(\pi\phi k)\,}.
\]
The sequence is deterministic, non-periodic, and bounded:
\[
\frac{\eta}{b}\le h_k\le \frac{\eta}{a}.
\]
By tuning $(a,b,\eta)$ to a target spectral interval $[\mu,L]$, one can place
$r_k/\eta$ inside a desired inverse-stepsize range (for instance within
$[\mu,L]$ in a quadratic model).

Compared with directly using Chebyshev roots in a fixed cycle, this construction avoids
explicit dependence on the horizon length $n$ and can be streamed online.

\section{Analysis}

\subsection{Polynomial viewpoint}

For the quadratic model, after $n$ iterations
\[
p_n(t)=\prod_{k=1}^n\left(1-\frac{\eta t}{r_k}\right),
\qquad
\|e_n\|_2\le \max_{\lambda\in[\mu,L]}|p_n(\lambda)|\,\|e_0\|_2.
\]
Hence performance is determined by how roots $\{r_k/\eta\}$ sample $[\mu,L]$ and by
their ordering. The golden-ratio rotation
$k\mapsto \{\phi k\}$ is a standard low-discrepancy mechanism, suggesting improved
spectral coverage over finite prefixes.

\subsection{Empirical boundedness phenomenon}

Define
\[
q_n(t)=\prod_{k=1}^n\left(1-\frac{t}{\sin^2(\pi\phi k)}\right),\qquad q_n(0)=1.
\]
Numerically, $q_n$ appears uniformly bounded on $t\in[0,1]$ over wide ranges of $n$,
and this behavior is highly sensitive to replacing $\phi$ with nearby irrational or
rational values.

\begin{conjecture}[Golden-root boundedness]
There exists a constant $C>0$ such that
\[
\sup_{n\ge 1}\sup_{t\in[0,1]}|q_n(t)|\le C.
\]
\end{conjecture}

If true, this would provide a direct path to non-asymptotic contraction bounds for
golden schedules after affine rescaling to a general spectral interval.

\begin{proposition}[Basic stability condition]
Let $A=A^\top\succ 0$ with spectrum in $[\mu,L]$, and suppose
$0<h_k<2/L$ for all $k$. Then gradient descent with variable steps satisfies
\[
\|e_n\|_2\le \prod_{k=1}^n \max\{|1-h_k\mu|,\ |1-h_kL|\}\,\|e_0\|_2.
\]
\end{proposition}

The proposition is standard and gives a conservative condition. The practical interest
here is precisely in schedules that may transiently exceed monotone-descent regimes but
still produce faster net contraction over blocks of iterations.

\section{Numerical tests}

\subsection{Quadratic optimisation: a solver for ill-conditioned linear systems}

We compare:
\begin{itemize}
\item constant-stepsize gradient descent,
\item Chebyshev-cyclic schedule,
\item golden schedule $h_k=\eta/(a+(b-a)\sin^2(\pi\phi k))$.
\end{itemize}

Test problem: diagonal SPD systems with prescribed condition number
$\kappa=L/\mu$ and random orthogonal similarity transforms. We report
$f(x_k)-f(x_\star)$ and $\|x_k-x_\star\|_2$ versus $k$.

Primary claim to test: for large $\kappa$, the golden schedule achieves a slope
consistent with accelerated behavior (empirically near $O(1/n^2)$ in objective gap)
without momentum terms.

\subsection{Neural network: handwritten digit recognition on MNIST}

We use the golden schedule as a learning-rate rule inside SGD/mini-batch GD:
\[
\alpha_k=\frac{\alpha_0}{a+(b-a)\sin^2(\pi\phi k)}.
\]
Baseline comparisons:
\begin{itemize}
\item constant learning rate,
\item cosine annealing,
\item one-cycle schedule.
\end{itemize}

Metrics: training loss, validation accuracy, robustness to $(a,b,\alpha_0)$, and
sensitivity to replacing $\phi$ by nearby constants.

Implementation note: experiments can be run in either PyTorch or TensorFlow; the
schedule only requires a one-line update rule at each optimizer step.

\subsection{Partial differential equations: the obstacle problem}

Discretizing the elliptic obstacle problem leads to large sparse convex quadratic (or
piecewise quadratic) subproblems. We evaluate whether golden long-step schedules reduce
iteration counts in projected-gradient style solvers while preserving feasibility.

We report:
\begin{itemize}
\item objective residual,
\item complementarity/constraint violation,
\item wall-clock time and iteration complexity.
\end{itemize}

\paragraph{Current status.}
This draft establishes the model, schedule definition, and testable conjecture.
Next steps are to formalize the boundedness proof (or a weaker finite-$n$ bound),
run the three experiment blocks, and refine the literature positioning with verified
citations.

\end{document}
